% Part: first-order-logic
% Chapter: natural-deduction
% Section: proving-things

\documentclass[../../../include/open-logic-section]{subfiles}

\begin{document}

\iftag{FOL}
      {\olfileid{fol}{ntd}{pro}}
      {\olfileid{pl}{ntd}{pro}}

\olsection{推导示例}

\begin{ex}
我们来给出语句 $(!A \land !B) \lif !A$ 的一个推导。

我们首先将待证的结论写在推导的底部。
\begin{prooftree}
\AxiomC{}
\UnaryInfC{$(!A\land !B) \lif !A$}
\end{prooftree}

接下来，我们需要弄清楚何种推理能得到这种形式的语句。结论的主联结词是 $\lif$，因此我们将尝试使用 \Intro{\lif} 规则来得到结论。在推导过程中，最好将涉及的假设写下来并标明所用的推理规则，以便在证明结束时容易看出所有假设是否均已解除。
\begin{prooftree}
\AxiomC{$\Discharge{!A \land !B}{1}$}
\DeduceC{$!A$}
\DischargeRule{\Intro{\lif}}{1} 
\UnaryInfC{$(!A\land !B) \lif !A$}
\end{prooftree}

现在我们需要补全从假设 $!A \land !B$ 到 $!A$ 的步骤。由于我们只需处理一个联结词 $\land$，因此必须使用 $\land$ 消去规则。这就得到了如下推导：
\begin{prooftree}
\AxiomC{$\Discharge{!A \land !B}{1}$}
\RightLabel{\Elim{\land}}
\UnaryInfC{$!A$}
\DischargeRule{\Intro{\lif}}{1} 
\UnaryInfC{$(!A\land !B) \lif !A$}
\end{prooftree}
这样，我们就得到了 $(!A \land !B) \lif !A$ 的一个正确推导。
\end{ex}

\begin{ex}
现在，让我们来给出 $(\lnot !A \lor !B)
\lif (!A \lif !B)$ 的一个推导。

我们首先将待证的结论写在推导的底部。
\begin{prooftree}
\AxiomC{}
\UnaryInfC{$(\lnot !A \lor !B) \lif (!A \lif !B)$}
\end{prooftree}
为了找到能得到该结论的逻辑规则，我们考察结论中的逻辑联结词：$\lnot$、$\lor$ 和 $\lif$。我们目前只关心第一个出现的 $\lif$，因为它是最终相继式中语句的主联结词，而 $\lnot$、$\lor$ 以及第二个出现的 $\lif$ 都在另一个联结词的辖域内，所以我们将稍后再处理它们。因此，我们从 \Intro{\lif} 规则开始。正确的应用必须形如：
\begin{prooftree}
\AxiomC{$\Discharge{\lnot !A \lor !B}{1}$}
\DeduceC{$!A \lif !B$}
\DischargeRule{\Intro{\lif}}{1}
\UnaryInfC{$(\lnot !A \lor !B) \lif (!A \lif !B)$}
\end{prooftree}
这给我们留下了两种继续的可能性。我们可以继续自底向上工作，寻找另一个 \Intro{\lif} 规则的应用，或者自顶向下工作，应用一个 \Elim{\lor} 规则。我们采用后者。我们将使用假设 $\lnot !A \lor !B$ 作为 \Elim{\lor} 规则最左边的前提。为了正确应用 \Elim{\lor} 规则，另外两个前提必须与结论 $!A \lif !B$ 相同，但它们每个都可以分别从另一个假设推导出来，即从 $\lnot !A \lor !B$ 的两个析取支之一推导出来。所以我们的推导将形如：
\begin{prooftree}
\AxiomC{$\Discharge{\lnot !A \lor !B}{1}$}
\AxiomC{$\Discharge{\lnot !A}{2}$}
\DeduceC{$!A \lif !B$}
\AxiomC{$\Discharge{!B}{2}$}
\DeduceC{$!A \lif !B$}
\DischargeRule{\Elim{\lor}}{2}
\TrinaryInfC{$!A \lif !B$}
\DischargeRule{\Intro{\lif}}{1} 
\UnaryInfC{$(\lnot !A \lor !B) \lif (!A \lif !B)$}
\end{prooftree}

在右边两个分支的每一个中，我们都需要推导 $!A \lif !B$，这最好使用 \Intro{\lif} 规则来完成。
\begin{prooftree}
\AxiomC{$\Discharge{\lnot !A \lor !B}{1}$}
\AxiomC{$\Discharge{\lnot !A}{2}, \Discharge{!A}{3}$}
\DeduceC{$!B$}
\DischargeRule{\Intro{\lif}}{3}
\UnaryInfC{$!A \lif !B$}
\AxiomC{$\Discharge{!B}{2}, \Discharge{!A}{4}$}
\DeduceC{$!B$}
\DischargeRule{\Intro{\lif}}{4}
\UnaryInfC{$!A \lif !B$}
\DischargeRule{\Elim{\lor}}{2}
\TrinaryInfC{$!A \lif !B$}
\DischargeRule{\Intro{\lif}}{1} 
\UnaryInfC{$(\lnot !A \lor !B) \lif (!A \lif !B)$}
\end{prooftree}

对于推导中缺失的两个部分，我们需要补全：从（中间分支的）$\lnot !A$ 和 $!A$ 推出 $!B$ 的推导，以及从（左边分支的）$!A$ 和 $!B$ 推出 $!B$ 的推导。我们先来处理前者。$\lnot !A$ 和 $!A$ 是 \Elim{\lnot} 规则的两个前提：
\begin{prooftree}
\AxiomC{$\Discharge{\lnot !A}{2}$}
\AxiomC{$\Discharge{!A}{3}$}
\RightLabel{\Elim{\lnot}}
\BinaryInfC{$\lfalse$}
\DeduceC{$!B$}
\end{prooftree}
通过使用 \FalseInt{} 规则，我们可以得到结论 $!B$，从而完成该分支。
\begin{prooftree}
\AxiomC{$\Discharge{\lnot !A \lor !B}{1}$}
\AxiomC{$\Discharge{\lnot !A}{2}$}
\AxiomC{$\Discharge{!A}{3}$}
\RightLabel{\Intro{\lfalse}}
\BinaryInfC{$\lfalse$}
\RightLabel{\FalseInt}
\UnaryInfC{$!B$}
\DischargeRule{\Intro{\lif}}{3}
\UnaryInfC{$!A \lif !B$}
\AxiomC{$\Discharge{!B}{2}, \Discharge{!A}{4}$}
\DeduceC{$!B$}
\DischargeRule{\Intro{\lif}}{4}
\UnaryInfC{$!A \lif !B$}
\DischargeRule{\Elim{\lor}}{2}
\TrinaryInfC{$!A \lif !B$}
\DischargeRule{\Intro{\lif}}{1} 
\UnaryInfC{$(\lnot !A \lor !B) \lif (!A \lif !B)$}
\end{prooftree}

现在来看最右边的分支。这里很重要的一点是，要意识到推导的定义\emph{允许解除假设}，但\emph{并不要求}必须解除。换句话说，如果我们能从假设 $!A$ 或 $!B$ 之一推出 $!B$ 而不使用另一个，也是可以的。而从~$!B$ 推出~$!B$ 是平凡的：$!B$ 本身就是一个这样的推导，不需要任何推理步骤。因此我们可以直接删去假设~$!A$。
\begin{prooftree}
\AxiomC{$\Discharge{\lnot !A \lor !B}{1}$}
\AxiomC{$\Discharge{\lnot !A}{2}$}
\AxiomC{$\Discharge{!A}{3}$}
\RightLabel{\Elim{\lnot}}
\BinaryInfC{$\lfalse$}
\RightLabel{\FalseInt}
\UnaryInfC{$!B$}
\DischargeRule{\Intro{\lif}}{3}
\UnaryInfC{$!A \lif !B$}
\AxiomC{$\Discharge{!B}{2}$}
\RightLabel{\Intro{\lif}}
\UnaryInfC{$!A \lif !B$}
\DischargeRule{\Elim{\lor}}{2}
\TrinaryInfC{$!A \lif !B$}
\DischargeRule{\Intro{\lif}}{1} 
\UnaryInfC{$(\lnot !A \lor !B) \lif (!A \lif !B)$}
\end{prooftree}
注意，在最终的推导中，最右边的 \Intro{\lif} 推理实际上并没有解除任何假设。
\end{ex}

\begin{ex}
到目前为止，我们还没有用到 \FalseCl{} 规则。它的特殊之处在于，它允许我们解除一个并非该规则结论之子公式的假设。它与 \FalseInt{} 规则密切相关。事实上，\FalseInt{} 规则是 \FalseCl{} 规则的一个特例——有一种称为“直觉主义逻辑”的逻辑，其中只允许使用 \FalseInt{} 规则。\FalseCl{} 规则是当其他方法都无效时的最后手段。例如，假设我们想推导 $!A \lor \lnot !A$。我们通常的策略会是尝试用 $\Intro{\lor}$ 来推导 $!A \lor \lnot !A$。但这会要求我们在没有任何假设的情况下推导出 $!A$ 或 $\lnot !A$，而这无法做到。这时就要靠 \FalseCl{} 了！
\begin{prooftree}
  \AxiomC{$\Discharge{\lnot(!A \lor \lnot !A)}{1}$}
  \DeduceC{$\lfalse$}
  \DischargeRule{\FalseCl}{1}
  \UnaryInfC{$!A \lor \lnot !A$}
\end{prooftree}
现在我们需要寻找一个从 $\lnot(!A \lor \lnot !A)$ 到 $\lfalse$ 的推导。由于 $\lfalse$ 是 \Elim{\lnot} 规则的结论，我们可以尝试应用该规则：
\begin{prooftree}
  \AxiomC{$\Discharge{\lnot(!A \lor \lnot !A)}{1}$}
  \DeduceC{$\lnot !A$}
  \AxiomC{$\Discharge{\lnot(!A \lor \lnot !A)}{1}$}
  \DeduceC{$!A$}
  \RightLabel{\Elim{\lnot}}
  \BinaryInfC{$\lfalse$}
  \DischargeRule{\FalseCl}{1}
  \UnaryInfC{$!A \lor \lnot !A$}
\end{prooftree}
为了寻找 $\lnot !A$ 的推导，我们的策略是应用 $\Intro{\lnot}$ 规则：
\begin{prooftree}
  \AxiomC{$\Discharge{\lnot(!A \lor \lnot !A)}{1}, \Discharge{!A}{2}$}
  \DeduceC{$\lfalse$}
  \DischargeRule{\Intro{\lnot}}{2}
  \UnaryInfC{$\lnot !A$}
  \AxiomC{$\Discharge{\lnot(!A \lor \lnot !A)}{1}$}
  \DeduceC{$!A$}
  \RightLabel{\Elim{\lnot}}
  \BinaryInfC{$\lfalse$}
  \DischargeRule{\FalseCl}{1}
  \UnaryInfC{$!A \lor \lnot !A$}
\end{prooftree}
这里，我们可以通过对假设 $\lnot(!A \lor \lnot !A)$ 以及由新假设 $!A$ 通过 $\Intro{\lor}$ 规则推出的 $!A \lor \lnot !A$ 应用 \Elim{\lnot} 规则，来轻松得到 $\lfalse$：
\begin{prooftree}
  \AxiomC{$\Discharge{\lnot(!A \lor \lnot !A)}{1}$}
  \AxiomC{$\Discharge{!A}{2}$}
  \RightLabel{\Intro{\lor}}
  \UnaryInfC{$!A \lor \lnot !A$}
  \RightLabel{\Elim{\lnot}}
  \BinaryInfC{$\lfalse$}
  \DischargeRule{\Intro{\lnot}}{2}
  \UnaryInfC{$\lnot !A$}
  \AxiomC{$\Discharge{\lnot(!A \lor \lnot !A)}{1}$}
  \DeduceC{$!A$}
  \RightLabel{\Elim{\lnot}}
  \BinaryInfC{$\lfalse$}
  \DischargeRule{\FalseCl}{1}
  \UnaryInfC{$!A \lor \lnot !A$}
\end{prooftree}
在右侧我们使用了相同的策略，只不过我们通过 \FalseCl{} 规则得到 $!A$：
\begin{prooftree}
  \AxiomC{$\Discharge{\lnot(!A \lor \lnot !A)}{1}$}
  \AxiomC{$\Discharge{!A}{2}$}
  \RightLabel{\Intro{\lor}}
  \UnaryInfC{$!A \lor \lnot !A$}
  \RightLabel{\Elim{\lnot}}
  \BinaryInfC{$\lfalse$}
  \DischargeRule{\Intro{\lnot}}{2}
  \UnaryInfC{$\lnot !A$}
  \AxiomC{$\Discharge{\lnot(!A \lor \lnot !A)}{1}$}
  \AxiomC{$\Discharge{\lnot !A}{3}$}
  \RightLabel{\Intro{\lor}}
  \UnaryInfC{$!A \lor \lnot !A$}
  \RightLabel{\Elim{\lnot}}
  \BinaryInfC{$\lfalse$}
  \DischargeRule{\FalseCl}{3}
  \UnaryInfC{$!A$}
  \RightLabel{\Elim{\lnot}}
  \BinaryInfC{$\lfalse$}
  \DischargeRule{\FalseCl}{1}
  \UnaryInfC{$!A \lor \lnot !A$}
\end{prooftree}
\end{ex}

\begin{prob}
给出推导以证明如下各式：
\begin{enumerate}
\item $!A \land (!B \land !C) \Proves (!A \land !B) \land !C$.
\item $!A \lor (!B \lor !C) \Proves (!A \lor !B) \lor !C$.
\item $!A \lif (!B \lif !C) \Proves !B \lif (!A \lif !C)$.
\item $!A \Proves \lnot\lnot !A$.
\end{enumerate}
\end{prob}

\begin{prob}
给出推导以证明如下各式：
\begin{enumerate}
\item $(!A \lor !B) \lif !C \Proves !A \lif !C$.
\item $(!A \lif !C) \land (!B \lif !C) \Proves (!A \lor !B) \lif !C$.
\item $\Proves \lnot(!A \land \lnot !A)$.
\item $!B \lif !A \Proves \lnot !A \lif \lnot !B$.
\item $\Proves (!A \lif \lnot !A) \lif \lnot !A$.
\item $\Proves \lnot(!A \lif !B) \lif \lnot !B$.
\item $!A \lif !C \Proves \lnot (!A \land \lnot !C)$.
\item $!A \land \lnot !C \Proves \lnot (!A \lif !C)$.
\item $!A \lor !B, \lnot !B \Proves !A$.
\item $\lnot !A \lor \lnot !B \Proves \lnot(!A \land !B)$.
\item $\Proves (\lnot !A \land \lnot !B) \lif\lnot(!A \lor !B)$.
\item $\Proves \lnot(!A \lor !B) \lif (\lnot !A \land \lnot !B)$.
\end{enumerate}
\end{prob}

\begin{prob}
给出推导以证明如下各式：
\begin{enumerate}
\item $\lnot(!A \lif !B) \Proves !A$.
\item $\lnot(!A \land !B) \Proves \lnot !A \lor \lnot !B$.
\item $!A \lif !B \Proves \lnot !A \lor !B$.
\item $\Proves \lnot \lnot !A \lif !A$.
\item $!A \lif !B, \lnot !A \lif !B \Proves !B$.
\item $(!A \land !B) \lif !C \Proves (!A \lif !C) \lor (!B \lif !C)$.
\item $(!A \lif !B) \lif !A \Proves !A$.
\item $\Proves (!A \lif !B) \lor (!B \lif !C)$.
\end{enumerate}
(这些题目都需要用到 $\FalseCl$~规则。)
\end{prob}

\end{document}